\chapter{Dermabrasion}

\section*{Overview}
Dermabrasion is a skin-resurfacing procedure in which the outer layers of damaged skin are mechanically abraded to improve its appearance and texture. It is used to treat scars, including acne scars, and to reduce the visibility of wrinkles, sun damage, and tattoos.
\section*{Alternative Names}
Surgical skin planing; cosmetic skin resurfacing; microdermabrasion (superficial variant)
\section*{Principle}
A rapidly rotating abrasive instrument removes the epidermis and superficial dermis, stimulating new collagen formation and re-epithelialization from adnexal structures in the deeper skin. Controlled depth of abrasion determines how aggressively the skin remodels while protecting deeper healing layers.
\section*{History}
Dermabrasion was developed in the early 20th century for scar revision and gained popularity after World War II for treating acne scarring. Modern variants include microdermabrasion and laser resurfacing.
\section*{Why It Is Done}
Ordered for acne scars, surgical and traumatic scars, facial wrinkles, actinic (sun) damage, and removal of superficial tattoos, usually when less invasive treatments have failed.
\section*{Is This a Primary or Secondary Test or Procedure}
Primary therapeutic procedure for selected skin textural and scarring disorders.
\section*{How It Is Performed}
The skin is cleaned, and a topical or injected anesthetic is applied, often with skin chilling or freezing. A handheld motorized abrader, such as a diamond burr or wire brush, is passed over the treatment area to a controlled depth. The wound is covered with a dressing or treated with an open-wound healing protocol.
\section*{Preparation}
Avoidance of isotretinoin for at least six months before the procedure because of abnormal wound healing, and pretreatment with antiviral medication when treating the perioral region to prevent herpes simplex flares.
\section*{Contraindications and Precautions}
Contraindications include active acne, skin infections, herpes simplex outbreak, and a history of keloid or hypertrophic scarring. Darker skin types have a higher risk of post-inflammatory pigmentary change and require extra caution.
\section*{Risks}
Pain, swelling, redness, infection, scarring, hyperpigmentation or hypopigmentation, milia, and prolonged healing with deeper abrasion.
\section*{Aftercare and Recovery}
The treated area crusts and heals over one to two weeks, with redness persisting for several weeks. Sun protection and gentle skin care are required during healing.
\section*{Results and Interpretation}
Improvement in scar depth, skin texture, and wrinkle appearance is gradual over weeks to months as collagen remodeling continues. Persistent or severe textural defects may require repeat treatment.
\section*{Expected Results}
Evenly resurfaced skin with a softer scar appearance and no significant pigmentary abnormality or scarring.
\section*{Follow-up}
Patients are followed for wound healing and pigmentary change, and repeat sessions may be scheduled for optimal results. Procedures for further resurfacing include microdermabrasion and laser therapy.
\section*{References}
\begin{enumerate}
  \item MedlinePlus. Dermabrasion. U.S. National Library of Medicine; 2025.
  \item Current Medical Diagnosis and Treatment. McGraw-Hill; 2025.
\end{enumerate}

\clearpage
